\begin{definition}
Fix the graph, palette, distinguished edge, exact terminal tree \(T\),
generating coloring \(\varphi_n\), positive-rung ETT, hierarchy base,
minimum hierarchy counterexample \(C\), and its
\noderef{HierarchyMinimumCounterexampleContext}.  A
\emph{Situation 7.2 tail datum} consists of the following exact data.
Put \(\ell=\noderef{HierarchyFullTailLength}(C)\).  One has
\(\ell\geq2\) and
\[
 \noderef{HierarchyPathNumber}(C.\mathit{hierarchy})=\ell .
\]
There are indices \(s\in\operatorname{Fin}|T|\) and
\(a,b\in\operatorname{Fin}(|T|+1)\) satisfying
\[
 s+1=|T|,\qquad a+1=|T|,\qquad b=|T|.
\]
Writing \(e_p=T.\mathit{edge}(s)\),
\(y_{p-1}=T.\mathit{vertex}(a)\), and
\(y_p=T.\mathit{vertex}(b)\), the exact edge \(e_p\) is not incident
with \(y_{p-1}\).  Finally a color \(\theta\) is selected with
\[
          \varphi_n(e_p)=\theta,\qquad
          \theta\notin\overline\varphi_n(y_{p-1}).
\]
This is precisely the common tail and Subcase .1 data used by all
three cases of Situation 7.2; parallel edges remain distinct.
\end{definition}
